\clearpage\begingroup\SourceLeaf{119b}\HistoricalHead{114}{L}{PHYSIK. BRIEFWECHSEL.}
\begin{center}[29.]\end{center}
\begin{center}\textsc{Gauss} an \textsc{Herger}.\end{center}
\noindent\hspace*{.18\linewidth}Wohlgeborener\par
\noindent\hspace*{.25\linewidth}Hochzuehrender Herr.\par\medskip
\indent Für das gütige Geschenk, welches Sie mir mit der ersten Lieferung Ihres
schätzbaren auch in seinem Aeussern so elegant ausgestatteten Werkes ge-
\clearpage\SourceLeaf{120}\HistoricalHead{115}{R}{GAUSS AN HERGER.}
\noindent macht haben\textsuperscript{[*)]}, verfehle ich nicht, Ihnen meinen schuldigen verbindlichsten
Dank abzustatten. Der von Ihnen dabei bewiesene ausdauernde Fleiss, sowie
die Accuratesse und Nettigkeit der Ausführung gereichen Ihnen zur Ehre und
verdienen alle Anerkennung. Ihre Curven sind in der That nichts anderes
als die rechtwinklichten Trajectorien (denn so wird ohne Zweifel das unleserlich
geschriebene Wort in Herrn von \textsc{Humboldts} Briefe gelesen werden müssen)
zu denjenigen Curven, welche nach meiner Theorie gleiche Potentiale vorstellen.
Aber die Gestaltung des einen wie des andern Curvensystems in
\textls[120]{der Nähe} von künstlichen Magneten lässt sich theoretisch oder a priori nicht
bestimmen, ohne die Vertheilung des Magnetismus im Innern von diesen vorher
zu kennen, und so bieten, wie Hr. Pr[ofessor] \textsc{Ermann} richtig bemerkt, Ihre
Curven ein sehr schätzbares Material dar, wonach jede vorgebrachte oder vorzubringende
Hypothese über solche Vertheilung wird geprüft werden können.
Dürfte ich mir erlauben, noch einen Wunsch auszusprechen, so wäre es der,
dass in der Einleitung über die Versuche selbst, auf welche die gezeichneten
Curven sich gründen, ein vollständigeres Detail mitgetheilt würde, ein Wunsch,
den Sie vielleicht schon selbst im Fortgange der Einleitung zu erfüllen sich
vorgesetzt haben, da dies in der That das am meisten geeignete Mittel ist,
das Zutrauen zu Ihren Resultaten zu befestigen und zu erhöhen.

Genehmigen Sie die Bezeugung der besonderen Hochachtung womit ich
beharre Ewr. Wohlgeboren
\begin{flushright}ergebenster Diener\hspace*{.12\linewidth}\\[6pt]
C. F. \textsc{Gauss}.\hspace*{.16\linewidth}\end{flushright}
\indent Göttingen, 11. Februar 1845.
\begingroup\renewcommand{\thefootnote}{}\footnotetext{[*) J. \textsc{Ernst Herger}, \textit{Die Systeme der magnetischen Curven, Isogonen und Isodynamen, nebst anderweitigen empirischen Fortschritten über die magnetisch-polaren Kräfte, ausgeführt in 37 grossen graphischen Darstellungen auf 31 Tafeln und erläutert unter den Auspicien des Hofrath Dr. Schottin.} Nebst einem Vorwort von G. A. \textsc{Erman}. In 4 Heften, Leipzig 1845. — Der Verfasser dieses Werkes J. E. \textsc{Herger} (1812—1871) lebte als Rosenzüchter und Handelsgärtner zu Köstritz im Fürstentum Reuß.]}\endgroup
\clearpage\endgroup
